% !TEX root = ../main.tex
\columnbreak
% =====================================================================
\section{Transformatoren}
% =====================================================================
\begin{symtab}
U & Spannung & [V] \\
u$_1$ & Spannung Primärseite & [V] \\
u$_2$ & Spannung Sekundärseite & [V] \\
I & Strom & [A] \\
i$_1$ & strom Primärseite & [A] \\
i$_2$ & Strom Sekundärseite & [A] \\
i$_{m1}$ & magnetisierungsstrom Primärseite & [A] \\
H & Feldstärke & [A/m] \\
N & Anzahl Wicklungen & [] \\
N$_1$ & Wicklungen Primärseite & [] \\
N$_2$ & Wicklungen Sekundärseite & [] \\
$\phi$ & Magnetischer Fluss & [Vs] \\
ü & Übersetzungsverhältnis & [] \\
ü$_k$ & Spannungsfaktor damit Nennstrom fliesst & [] \\
R$_{fe}$ & Ummagnetisierungsverluste & [$\Omega$] \\
R$_1$ & Wicklungswiderstand Primärseite & [$\Omega$] \\
R$_2$ & Wicklungswiderstand Sekundärseite & [$\Omega$] \\
L$_{h1}$ & Streuinduktivität & [H] \\
L$_{\sigma 1}$ & Induktivität Primärseite & [H] \\
L$_{\sigma 1}$ & Induktivität Sekundärseite & [H] \\
\end{symtab}

\subsection{Ideal}
\fig[0.3]{trafo_ideal}
{\setlength{\jot}{3pt}
\begin{gather*}
U = N \cdot \frac{d\phi}{dt} \\
\phi = \frac{1}{N} \int U dt \\
\text{ü} = \frac{i_2}{i_1} = \frac{u_1}{u_2} = \frac{N_1}{N_2}
\end{gather*}}

\subsubsection{Magnetisierungsstrom}
\fig[0.6]{trafo_magnetisierung}
{\setlength{\jot}{3pt}
\begin{gather*}
L_{h1} = \mu_0 \cdot \mu_r \cdot \frac{A_{fe}}{l_{fe}} \cdot N_1^2 \\
i_{m1} = i_1 - \frac{N_2}{N_1} \cdot i_2
\end{gather*}}

\columnbreak
{\setlength{\jot}{3pt}
\begin{gather*}
L_{h2} = \frac{L_{h1}}{\text{ü}^2} \\
\phi = \mu_0 \cdot \mu_r \cdot A_{fe} \cdot H \\
u_1 = L_{h1} \cdot \frac{di_{m1}}{dt}
\end{gather*}}

\subsubsection{Streuung}
\fig[0.6]{trafo_streuung}

\subsection{Real}
\subsubsection{Wicklungsverluste}
\fig[0.75]{trafo_wicklungsverluste}
{\setlength{\jot}{3pt}
\begin{gather*}
u_1 = N_1 \cdot \frac{d}{dt}(\phi_h + \phi_{\sigma 1}) \\
u_2 = N_2 \cdot \frac{d}{dt}(\phi_h + \phi_{\sigma 1}) % TODO prüfen: sekundär wohl phi_sigma2
\end{gather*}}

\subsubsection{Bezug auf Primärseite}
\fig[0.9]{trafo_bezug_primaer}

\columnbreak
\fig[0.9]{trafo_t_ersatzschaltung}
{\setlength{\jot}{3pt}
\begin{gather*}
u'_2 = u_2 \text{ü} \\
i'_2 = \frac{i_2}{\text{ü}} \\
R'_2 = R_2 \text{ü}^2 \\
L'_{\sigma 2} = L_{\sigma 2} \text{ü}^2
\end{gather*}}

\subsubsection{Leerlaufversuch}
Bei den Versuchen werden komplexe Wechselstromrechnungen mit $\|Z\|$ gemacht.
\fig[0.9]{trafo_leerlauf}

\subsubsection{Kurzschlussversuch}
\fig[0.63]{trafo_kurzschluss}
ü$_k$ wird ofmals als u$_k$ in \% angegeben\par
Oft wird folgende Annahme gemacht:
\begin{gather*}
L_{\sigma 1} = L'_{\sigma 2} = \frac{L_{\sigma 1} + L'_{\sigma 2}}{2}
\end{gather*}
R$_1$ und R$_2$ lassen sich direkt messen.\par
Spetzialfall falls L$_\sigma$ nicht gegeben und L$_h$ gesucht ist.
\begin{gather*}
\omega L_h = \sqrt{Z_{Kurzschluss}^2 - R_{Gemessengesammt}^2}
\end{gather*}

\columnbreak
\subsubsection{Sättigung}
\fig[0.7]{trafo_saettigung}
